\scanpage{6}
\begin{center}
IV. \uline{The Brauer group of a ringed space.}
\end{center}

The foregoing theorems enable us to define the Brauer group of a ringed space $(\mathrm{X},\handwritten{\mathscr{A}})$, where $\handwritten{\mathscr{A}}$ is a sheaf of commutative rings with unit.

We consider the category $\handwritten{\mathscr{C}}_1(\mathrm{X},\handwritten{\mathscr{A}})$ of locally separable $\handwritten{\mathscr{A}}$-algebras, with arrows the $\handwritten{\mathscr{A}}$-algebra morphisms. $\handwritten{\mathscr{C}}_1(\mathrm{X},\handwritten{\mathscr{A}})$ is closed under sheaf tensor product by Theorem III.5. We define $\handwritten{\mathscr{C}}'_1(\mathrm{X},\handwritten{\mathscr{A}})$ to be the full subcategory of $\handwritten{\mathscr{C}}_1(\mathrm{X},\handwritten{\mathscr{A}})$ consisting of those objects which are isomorphic to the algebras of the form $\underline{\mathrm{Hom}}_{\handwritten{\mathscr{A}}}(\handwritten{\mathscr{M}},\handwritten{\mathscr{M}})$, where $\handwritten{\mathscr{M}}$ is a faithful locally projective $\handwritten{\mathscr{A}}$-module of finite type (Theorem III,4). From Theorem II.4, we know that $\handwritten{\mathscr{C}}'_1(\mathrm{X},\handwritten{\mathscr{A}})$ is closed under tensor product.

We define an equivalence relation in $\handwritten{\mathscr{C}}_1(\mathrm{X},\handwritten{\mathscr{A}})$ as follows : $\mathscr{L}_1$ and $\mathscr{L}_2$ in $\handwritten{\mathscr{C}}_1(\mathrm{X},\handwritten{\mathscr{A}})$ are equivalent if and only if there exist $\handwritten{\mathcal{M}}_1$ and $\handwritten{\mathcal{M}}_2$ in $\handwritten{\mathscr{C}}'_1(\mathrm{X},\handwritten{\mathscr{A}})$ such that $\mathscr{L}_1\otimes_{\handwritten{\mathscr{A}}}\handwritten{\mathcal{M}}_1\approx \mathscr{L}_2\otimes_{\handwritten{\mathscr{A}}}\handwritten{\mathcal{M}}_2$.

The set of equivalence classes, $\handwritten{\mathfrak{B}}(\mathrm X,\handwritten{\mathscr{A}})$ is an abelian group under tensor product, with unit $\{\handwritten{\mathscr{A}}\}$ the class of $\handwritten{\mathscr{A}}$. If $\mathscr{L}$ is in $\handwritten{\mathscr{C}}_1(\mathrm{X},\handwritten{\mathscr{A}})$, then the inverse of $\{\mathscr{L}\}$ is $\{\mathscr{L}^{\unsure{\mathrm{o}|0}}\}$, since $\eta\colon \mathscr{L}\otimes_{\handwritten{\mathscr{A}}}\mathscr{L}^{\mathrm{o}}\to \underline{\mathrm{Hom}}_{\handwritten{\mathscr{A}}}(\mathscr{L},\mathscr{L})$ is an isomorphism.

THEOREM IV.1. If $\mathscr{L}$ is in $\handwritten{\mathscr{C}}_1(\mathrm{X},\handwritten{\mathscr{A}})$, then $\mathscr{L}$ is in $\handwritten{\mathscr{C}}'_1(\mathrm{X},\handwritten{\mathscr{A}})$
% archive sheet number: 5
